\section{True original dual transposes from actual canonical evaluation}

\begin{definition}[Actual transpose to the original sheaf dual]
\label{def:actual-invertible-dual-transpose}
\lean{AlgebraicGeometry.Scheme.Modules.originalRightUnitIso}
\lean{AlgebraicGeometry.Scheme.Modules.IsInvertible.dualTranspose}
\leanok
\uses{lem:cartier-actual-canonical-dual-evaluation-isomorphism,lem:cartier-actual-line-tensor-fully-faithful}
For an actual original invertible module $M$ and an actual whole
original morphism $\alpha:M\to\mathcal O_X$, define its transpose
$\sigma_\alpha:\mathcal O_X\to M^\vee$ by the actual full
faithfulness of tensoring with $M$. Its tensor image is the
actual original right unit isomorphism followed by $\alpha$ and
the inverse actual canonical dual evaluation. Thus the source,
target, and pairing are the original modules and original
canonical sheaf evaluation, with no abstract duality input.
\end{definition}

\begin{lemma}[Actual whole transpose pairing identity and uniqueness]
\label{lem:actual-invertible-dual-transpose-pairing}
\lean{AlgebraicGeometry.Scheme.Modules.IsInvertible.dualTranspose_pairing}
\lean{AlgebraicGeometry.Scheme.Modules.IsInvertible.dualTranspose_unique}
\leanok
\uses{def:actual-invertible-dual-transpose}
Tensoring the actual transpose with the original identity of $M$
and following the original canonical evaluation gives exactly
$\alpha$ after the true original right-unit comparison. This
whole identity uniquely determines the original transpose among
all original module morphisms.
\end{lemma}
\begin{proof}
Actual tensor full faithfulness sends the constructed preimage to
its stated whole tensor image. Postcompose by actual evaluation;
its genuine inverse-hom identity cancels the inverse evaluation
factor and gives the asserted whole original pairing equality.
For any other original morphism satisfying it, cancel the genuine
original evaluation isomorphism and then apply injectivity of the
same actual tensor functor on original hom-sets. The two entire
original morphisms are equal.
\end{proof}
